% =============================================================================
% Strength Labs — Documento de arquitectura (detallado)
% =============================================================================
\documentclass[11pt,a4paper]{article}

\usepackage[T1]{fontenc}
\usepackage[utf8]{inputenc}
\usepackage[spanish,es-noshorthands]{babel}
\usepackage{geometry}
\geometry{margin=2.4cm}
\usepackage{hyperref}
\usepackage{xcolor}
\usepackage{listings}
\usepackage{enumitem}
\usepackage{tabularx}
\usepackage{booktabs}
\usepackage{longtable}
\usepackage{array}
\usepackage{tikz}
\usepackage{float}
\usepackage{microtype}
\usepackage{titlesec}
\usepackage{fancyhdr}
\usetikzlibrary{positioning,arrows.meta,shapes.geometric,fit,backgrounds,calc}

\hypersetup{
  colorlinks=true,
  linkcolor=blue!50!black,
  urlcolor=blue!50!black,
  citecolor=blue!50!black,
}

% ── Estilo de listings ───────────────────────────────────────────────────────
\definecolor{lstbg}{RGB}{248,248,248}
\definecolor{lstkw}{RGB}{0,80,160}
\definecolor{lststr}{RGB}{160,0,80}
\definecolor{lstcom}{RGB}{0,100,40}
\lstset{
  basicstyle=\ttfamily\footnotesize,
  backgroundcolor=\color{lstbg},
  keywordstyle=\color{lstkw}\bfseries,
  stringstyle=\color{lststr},
  commentstyle=\color{lstcom}\itshape,
  breaklines=true,
  frame=single,
  framesep=4pt,
  rulecolor=\color{black!30},
  columns=fullflexible,
  showstringspaces=false,
  numbers=none,
  captionpos=b,
}

% ── Encabezado / pie ─────────────────────────────────────────────────────────
\pagestyle{fancy}
\fancyhf{}
\renewcommand{\headrulewidth}{0.4pt}
\fancyhead[L]{\small Strength Labs}
\fancyhead[R]{\small Arquitectura técnica}
\fancyfoot[C]{\thepage}

% ── Comandos de utilidad ─────────────────────────────────────────────────────
\newcommand{\file}[1]{\texttt{\small #1}}
\newcommand{\cls}[1]{\texttt{\small #1}}
\newcommand{\http}[1]{\texttt{\small #1}}

\title{
  \vspace*{1.2cm}
  {\Huge\bfseries Strength Labs}\\[0.6cm]
  {\Large Arquitectura técnica detallada}\\[0.4cm]
  \rule{0.6\linewidth}{0.4pt}\\[0.4cm]
  {\large Mobile training tracker — Flutter, Spring Boot, FastAPI}
}
\author{Yeisen Kenneth}
\date{Mayo 2026 — versión 2}

\begin{document}
\maketitle
\thispagestyle{empty}
\vfill
\begin{center}
\small
Repositorios: \url{github.com/YeisenK/StrengthLabs} (cliente Flutter),
\url{github.com/YeisenK/StrengthLabsBackend} (backend Spring Boot + compute Python).\\
Endpoint público: \url{https://api.blocsa.com}.
\end{center}
\clearpage

\tableofcontents
\listoffigures
\listoftables
\clearpage

% =============================================================================
\section{Resumen ejecutivo}
% =============================================================================

Strength Labs es una aplicación móvil multiplataforma para el registro de entrenamientos de fuerza con seguimiento de carga interna (ATL, CTL, TSB, ACWR) y modelos de riesgo de lesión y sobreentrenamiento. El sistema se compone de tres servicios distribuidos:

\begin{itemize}[leftmargin=*]
  \item Un cliente Flutter offline-first con base de datos local SQLite (Drift), cola de sincronización con reintentos idempotentes, y soporte para Android e iOS.
  \item Un backend Java + Spring Boot 4 con persistencia en PostgreSQL gestionada por Flyway, autenticación JWT RS256, control de concurrencia optimista y filtro de correlación por request.
  \item Un microservicio Python + FastAPI sin estado que ejecuta los modelos numéricos de fatiga y riesgo, consumido por el backend mediante HTTP con timeouts y degradación controlada ante fallos.
\end{itemize}

El propósito de este documento es servir como referencia completa para mantenedores, integradores y revisores. Cubre el modelo de datos, los contratos de API, los flujos de autenticación y sincronización, la estrategia de pruebas, el despliegue real (incluyendo el túnel Cloudflare con dominio \texttt{api.blocsa.com}), las decisiones arquitectónicas y sus trade-offs.

% =============================================================================
\section{Glosario}
% =============================================================================

\begin{longtable}{p{0.22\textwidth} p{0.72\textwidth}}
\toprule
\textbf{Término} & \textbf{Definición} \\
\midrule
\endhead
\textbf{ATL} & \textit{Acute Training Load}. Promedio exponencial de la carga de entrenamiento de los últimos 7 días. \\
\textbf{CTL} & \textit{Chronic Training Load}. Promedio exponencial de los últimos 42 días. \\
\textbf{TSB} & \textit{Training Stress Balance} = CTL $-$ ATL. Cuanto más negativo, más fatiga acumulada. \\
\textbf{ACWR} & \textit{Acute:Chronic Workload Ratio}. ATL/CTL. Por encima de $\sim$1{,}5 sube el riesgo de lesión. \\
\textbf{Monotony} & Desviación estándar de la carga diaria normalizada. Alta monotonía + alta carga $\Rightarrow$ alto strain. \\
\textbf{Strain} & Producto carga $\times$ monotony. Indicador agregado de estrés semanal. \\
\textbf{Ramp rate} & Variación semanal de CTL. Por encima de cierto umbral aumenta el riesgo. \\
\textbf{RPE} & \textit{Rate of Perceived Exertion}, escala 1–10 que el usuario asigna a una serie. \\
\textbf{Cubit} & Variante simplificada de Bloc; almacén de estado emitido vía streams. \\
\textbf{Drift} & ORM SQLite para Dart/Flutter con codegen y streams reactivos. \\
\textbf{Flyway} & Migrador de bases de datos basado en archivos SQL versionados. \\
\textbf{Idempotency key} & Identificador único por intención de mutación, enviado por el cliente para deduplicar reintentos. En este sistema se llama \texttt{client\_request\_id}. \\
\textbf{Optimistic locking} & Control de concurrencia donde cada fila lleva una versión; los updates fallan si la versión esperada no coincide con la actual. \\
\textbf{N+1 query} & Patrón anti-rendimiento de ORM donde cargar una lista dispara una consulta por elemento para asociaciones diferidas. \\
\textbf{Quick tunnel} & Túnel Cloudflare efímero (\texttt{trycloudflare.com}) sin cuenta. \\
\textbf{Named tunnel} & Túnel Cloudflare persistente bajo cuenta, con dominio fijo y reinicio sin rotación. \\
\bottomrule
\end{longtable}

% =============================================================================
\section{Vista de contexto}
% =============================================================================

\subsection{Diagrama C4 — nivel 1}

\begin{figure}[H]
\centering
\begin{tikzpicture}[
  actor/.style={rectangle,draw,minimum height=1.1cm,minimum width=2.6cm,align=center,fill=gray!10},
  system/.style={rectangle,draw,thick,rounded corners,minimum height=1.4cm,minimum width=4.4cm,align=center,fill=blue!12},
  external/.style={rectangle,draw,dashed,minimum height=1.1cm,minimum width=3.0cm,align=center,fill=orange!10},
  arrow/.style={-Latex,thick,font=\footnotesize},
]
  \node[actor] (user) {Usuario\\\small (atleta amateur)};
  \node[system,right=3.2cm of user] (sys) {Strength Labs\\\small Mobile + Backend};
  \node[external,right=3.0cm of sys, yshift=0.8cm] (google) {Google Identity};
  \node[external,right=3.0cm of sys, yshift=-0.8cm] (cf) {Cloudflare\\Edge};

  \draw[arrow] (user) -- node[above]{registra workouts\\\small consulta fatiga} (sys);
  \draw[arrow] (sys) -- (google);
  \draw[arrow] (sys) -- node[below right,xshift=-4pt,yshift=2pt]{TLS} (cf);
\end{tikzpicture}
\caption{Vista de contexto. Strength Labs ofrece al atleta un registro de entrenamientos con métricas de carga y se apoya en Google para autenticación opcional y en Cloudflare para la terminación TLS pública.}
\end{figure}

\subsection{Actores}

\begin{itemize}[leftmargin=*]
  \item \textbf{Atleta}. Es el único actor humano. Crea su cuenta (correo o Google), registra sesiones, consulta su tendencia de fatiga, y exporta sus datos a XLSX/CSV.
  \item \textbf{Servicio de identidad de Google}. Verificación del \textit{ID token} OAuth durante el flujo de Google Sign-In. No es indispensable: el sistema funciona con email/password.
  \item \textbf{Red Cloudflare}. Termina TLS, expone el dominio público \texttt{api.blocsa.com} y reenvía a \texttt{localhost:8080} en el host del backend mediante un \textit{named tunnel}.
\end{itemize}

\subsection{Drivers arquitectónicos}

\begin{table}[H]
\centering
\begin{tabularx}{\textwidth}{lX}
\toprule
\textbf{Driver} & \textbf{Implicación} \\
\midrule
Funcionalidad offline & La app debe permitir crear/editar workouts sin red y reconciliar al volver online. Esto fuerza una BD local, una cola de sync y claves de idempotencia. \\
Seguridad razonable & JWT firmados asimétricamente, claves fuera del repositorio, rate-limiting en auth, no exponer trazas al cliente. \\
Trazabilidad & Cada request lleva un \texttt{X-Request-Id} echoed en respuesta y MDC para correlacionar logs cliente/servidor. \\
Coste operativo bajo & Despliegue en un solo host (Postgres + Spring + FastAPI + Cloudflared) con scripts simples (systemd). \\
Alcance académico & No se invierte en infraestructura de alto nivel (Redis distribuido, Kubernetes, Sentry); las decisiones documentan deliberadamente dónde la solución se queda corta para producción real. \\
\bottomrule
\end{tabularx}
\caption{Drivers arquitectónicos.}
\end{table}

% =============================================================================
\section{Vista de contenedores}
% =============================================================================

\subsection{Diagrama C4 — nivel 2}

\begin{figure}[H]
\centering
\begin{tikzpicture}[
  box/.style={rectangle,draw,rounded corners,minimum height=1.1cm,minimum width=3.4cm,align=center,font=\small},
  app/.style={box,fill=blue!10},
  svc/.style={box,fill=blue!15},
  ext/.style={box,fill=orange!10},
  store/.style={cylinder,draw,shape border rotate=90,aspect=0.25,minimum height=1.0cm,minimum width=1.6cm,align=center,font=\small,fill=yellow!15},
  arrow/.style={-Latex,thick,font=\footnotesize},
]
  \node[app] (flutter) {Flutter app\\\small iOS / Android};
  \node[store,below=0.8cm of flutter] (drift) {Drift\\\small SQLite local};
  \node[ext,right=2.8cm of flutter,yshift=1.4cm] (cf) {Cloudflare Tunnel\\\small api.blocsa.com};
  \node[svc,right=2.8cm of cf,yshift=-1.0cm] (api) {Spring Boot API\\\small :8080};
  \node[svc,right=1.4cm of api] (compute) {Compute engine\\\small FastAPI :8001};
  \node[store,below=0.8cm of api] (pg) {PostgreSQL 16};

  \draw[arrow] (flutter) -- node[above,yshift=4pt]{HTTPS} (cf);
  \draw[arrow] (cf) -- node[above,yshift=2pt]{HTTP local} (api);
  \draw[arrow] (api) -- node[above]{HTTP} (compute);
  \draw[arrow] (api) -- (pg);
  \draw[arrow,<->] (flutter) -- (drift);
\end{tikzpicture}
\caption{Vista de contenedores. El cliente Flutter mantiene una base local Drift y habla únicamente con \texttt{api.blocsa.com}. El borde Cloudflare reenvía a \texttt{localhost:8080} en el host del backend, que coordina Postgres y el microservicio Python.}
\end{figure}

\subsection{Comunicación entre contenedores}

\begin{itemize}[leftmargin=*]
  \item \textbf{Flutter $\rightarrow$ Cloudflare $\rightarrow$ Spring}: HTTPS / HTTP REST. Una sola dirección base (\texttt{ApiConstants.baseUrl}). El cliente persiste todas las mutaciones en Drift antes de salir a red.
  \item \textbf{Flutter $\leftrightarrow$ Drift}: bidireccional, no es ``almacenamiento de respaldo'': es la fuente de verdad para reads.
  \item \textbf{Spring $\rightarrow$ Compute Engine}: POST JSON contra \texttt{/compute/fatigue} y \texttt{/compute/risk}. Timeouts (2\,s connect, 5\,s read) y fallback a valores neutros.
  \item \textbf{Spring $\leftrightarrow$ PostgreSQL}: JDBC vía HikariCP. Esquema gobernado por Flyway.
\end{itemize}

% =============================================================================
\section{Tecnologías y dependencias}
% =============================================================================

\begin{table}[H]
\centering
\begin{tabularx}{\textwidth}{lXX}
\toprule
\textbf{Capa} & \textbf{Frontend} & \textbf{Backend} \\
\midrule
Lenguaje & Dart 3.10 & Java 21 \\
Framework & Flutter 3 & Spring Boot 4.0.3 \\
Estado & \texttt{flutter\_bloc} (Cubit) & — \\
Navegación & \texttt{go\_router} & — \\
HTTP & \texttt{dio} 5 + interceptores & Spring MVC + RestTemplate \\
Persistencia local & \texttt{drift} (SQLite) & — \\
Persistencia remota & — & Spring Data JPA + Hibernate 7 \\
Migraciones & — & Flyway (\texttt{spring-boot-flyway}) \\
Auth & JWT + \texttt{google\_sign\_in} & Spring Security + JJWT 0.12.6 RS256 \\
Almacenamiento de claves & \texttt{flutter\_secure\_storage} & Variables de entorno (PEM base64) \\
Conectividad & \texttt{connectivity\_plus} & — \\
Identidad & \texttt{uuid} v4 (\textit{client\_request\_id}, ids locales) & \texttt{UUID.randomUUID} para PK \\
Pruebas & \texttt{flutter\_test} + \texttt{bloc\_test} + \texttt{mocktail} & JUnit 5 + Testcontainers Postgres \\
\bottomrule
\end{tabularx}
\caption{Stack tecnológico.}
\end{table}

El microservicio Python utiliza FastAPI con Pydantic para los modelos de entrada/salida y NumPy / pandas para los cálculos de carga.

% =============================================================================
\section{Backend: Spring Boot API}
% =============================================================================

\subsection{Capas hexagonales}

El backend separa responsabilidades en cinco paquetes Java alineados con la arquitectura hexagonal:

\begin{description}[leftmargin=2cm,style=nextline]
  \item[\file{com.strengthlabs.api}] Punto de entrada Spring Boot y configuración (seguridad, i18n, \texttt{RestTemplate} con timeouts).
  \item[\file{com.strengthlabs.application}] Casos de uso (\cls{GetFatigueSummaryUseCase}), DTOs internos y puertos hacia servicios externos.
  \item[\file{com.strengthlabs.domain}] Entidades de dominio (\cls{User}, \cls{Routine}) y reglas puras sin dependencia de framework.
  \item[\file{com.strengthlabs.infrastructure}] Adaptadores: entidades JPA, repositorios Spring Data, adaptador HTTP para el compute engine, infraestructura de seguridad (\cls{JwtTokenProvider}, \cls{RbacFilter}, \cls{LoginRateLimitFilter}, \cls{RevokedTokenStore}) y filtros web (\cls{CorrelationIdFilter}).
  \item[\file{com.strengthlabs.presentation}] Controladores REST (\cls{AuthController}, \cls{WorkoutController}, \cls{FatigueController}, \cls{RoutineController}, \cls{ExerciseController}) y middleware (\cls{GlobalExceptionHandler}, \cls{LocalizedStatusException}).
\end{description}

\subsection{Modelo de datos}

\subsubsection{Diagrama entidad–relación}

\begin{figure}[H]
\centering
\begin{tikzpicture}[
  entity/.style={rectangle,draw,minimum width=2.6cm,minimum height=0.9cm,align=center,font=\small,fill=yellow!10},
  arrow/.style={-Latex,thick,font=\scriptsize},
]
  \node[entity] (user) {users};
  \node[entity,right=2.6cm of user] (workout) {workouts};
  \node[entity,right=2.6cm of workout] (we) {workout\_exercises};
  \node[entity,right=2.6cm of we] (sets) {workout\_sets};
  \node[entity,below=0.9cm of we] (ex) {exercises};
  \node[entity,below=0.9cm of user] (metrics) {training\_metrics};
  \node[entity,below=0.9cm of workout] (rout) {routines};
  \node[entity,below=0.9cm of sets] (rdays) {routine\_days};

  \draw[arrow] (user) -- node[above]{1:N} (workout);
  \draw[arrow] (workout) -- node[above]{1:N} (we);
  \draw[arrow] (we) -- node[above]{1:N} (sets);
  \draw[arrow] (we) -- node[right]{N:1} (ex);
  \draw[arrow] (user) -- node[right]{1:N} (metrics);
  \draw[arrow] (rout) -- node[right]{1:N} (rdays);
\end{tikzpicture}
\caption{Modelo entidad–relación simplificado del backend.}
\end{figure}

\subsubsection{DDL clave (tras V6)}

\begin{lstlisting}[language=SQL,caption={Definición de las tablas \texttt{users}, \texttt{workouts} y \texttt{training\_metrics}.}]
CREATE TABLE users (
    id            UUID         PRIMARY KEY,
    name          VARCHAR(255) NOT NULL,
    email         VARCHAR(255) NOT NULL UNIQUE,
    password_hash TEXT         NOT NULL,
    role          VARCHAR(20)  NOT NULL DEFAULT 'USER',
    active        BOOLEAN      NOT NULL DEFAULT TRUE
);

CREATE TABLE workouts (
    id                UUID         PRIMARY KEY,
    user_id           UUID         NOT NULL REFERENCES users (id) ON DELETE CASCADE,
    name              VARCHAR(255) NOT NULL,
    date              TIMESTAMPTZ  NOT NULL,
    duration_seconds  INT          NOT NULL DEFAULT 0,
    notes             TEXT,
    client_request_id UUID,
    version           BIGINT       NOT NULL DEFAULT 0
);
CREATE UNIQUE INDEX uq_workouts_user_client_request
    ON workouts (user_id, client_request_id)
    WHERE client_request_id IS NOT NULL;

CREATE TABLE training_metrics (
    id                       UUID            PRIMARY KEY,
    user_id                  UUID            NOT NULL REFERENCES users (id) ON DELETE CASCADE,
    metric_date              DATE            NOT NULL,
    atl                      DOUBLE PRECISION NOT NULL DEFAULT 0,
    ctl                      DOUBLE PRECISION NOT NULL DEFAULT 0,
    tsb                      DOUBLE PRECISION NOT NULL DEFAULT 0,
    acwr                     DOUBLE PRECISION NOT NULL DEFAULT 0,
    monotony                 DOUBLE PRECISION NOT NULL DEFAULT 0,
    strain                   DOUBLE PRECISION NOT NULL DEFAULT 0,
    ramp_rate                DOUBLE PRECISION NOT NULL DEFAULT 0,
    readiness_score          DOUBLE PRECISION NOT NULL DEFAULT 0,
    injury_risk_score        DOUBLE PRECISION NOT NULL DEFAULT 0,
    overtraining_risk_score  DOUBLE PRECISION NOT NULL DEFAULT 0,
    composite_risk_score     DOUBLE PRECISION NOT NULL DEFAULT 0,
    risk_level               VARCHAR(20)     NOT NULL DEFAULT 'low',
    dominant_factor          VARCHAR(100)    NOT NULL DEFAULT '',
    risk_flags_json          TEXT            NOT NULL DEFAULT '[]',
    recommendations_json     TEXT            NOT NULL DEFAULT '[]',
    weekly_volume_json       TEXT            NOT NULL DEFAULT '{}',
    trend_json               TEXT            NOT NULL DEFAULT '[]',
    compute_available        BOOLEAN         NOT NULL DEFAULT FALSE,
    created_at               TIMESTAMPTZ     NOT NULL DEFAULT NOW(),
    CONSTRAINT uq_training_metrics_user_date UNIQUE (user_id, metric_date)
);
\end{lstlisting}

\subsubsection{Migraciones Flyway}

\begin{table}[H]
\centering
\begin{tabularx}{\textwidth}{lX}
\toprule
\textbf{Versión} & \textbf{Contenido} \\
\midrule
V1 & Esquema inicial: \texttt{users}, \texttt{exercises}, \texttt{workouts}, \texttt{workout\_exercises}, \texttt{workout\_sets}, \texttt{training\_metrics}. \\
V2 & Catálogo de rutinas y días. \\
V3 & Columnas \texttt{name\_es} para i18n del catálogo. \\
V4 & Traducciones de catálogo (cargas iniciales). \\
V5 & \texttt{client\_request\_id} y \texttt{version} en \texttt{workouts}; índice único parcial. \\
V6 & Convierte \texttt{NUMERIC(p,s)} a \texttt{DOUBLE PRECISION} en \texttt{training\_metrics} y \texttt{workout\_sets} para que la entidad JPA \texttt{double} valide. \\
\bottomrule
\end{tabularx}
\caption{Resumen de migraciones Flyway aplicadas.}
\end{table}

\paragraph{Caso especial — baseline.} En despliegues anteriores el backend corrió con \texttt{ddl-auto=update}, por lo que Hibernate construyó el esquema bypaseando Flyway. Al añadir el módulo \texttt{spring-boot-flyway} de Spring Boot 4 (necesario porque la autoconfiguración se separó del paquete principal) y arrancar, Flyway detectó un esquema poblado sin tabla \texttt{flyway\_schema\_history}. La resolución documentada es adoptar el esquema existente como baseline V4 mediante:

\begin{lstlisting}[caption={Recetas de baseline para adoptar un esquema construido por Hibernate.}]
SPRING_FLYWAY_BASELINE_ON_MIGRATE=true
SPRING_FLYWAY_BASELINE_VERSION=4
SPRING_FLYWAY_BASELINE_DESCRIPTION=adopt-hibernate-built-schema-as-V4
\end{lstlisting}

V5 y V6 se aplican encima sin pérdida de datos.

\subsection{Endpoints REST}

\subsubsection{Resumen}

\begin{longtable}{lp{0.42\textwidth}p{0.36\textwidth}}
\toprule
\textbf{Método} & \textbf{Ruta} & \textbf{Notas} \\
\midrule
\endhead
POST   & \http{/auth/register}        & Crea usuario, devuelve par de tokens. \\
POST   & \http{/auth/login}           & Email + password. Rate-limited por IP. \\
POST   & \http{/auth/refresh}         & Rota ambos tokens; revoca el JTI anterior. \\
POST   & \http{/auth/logout}          & Revoca el refresh token. \\
POST   & \http{/auth/google}          & Verifica ID token de Google, emite JWT propio. \\
GET    & \http{/auth/me}              & Perfil autenticado. \\
GET    & \http{/.well-known/jwks.json}& Clave pública para verificadores externos. \\
GET    & \http{/workouts}             & Sin params: \texttt{\{items\}}. Con \texttt{?page=\&size=}: \texttt{\{items, page, size, total, has\_more\}}. \\
POST   & \http{/workouts}             & Acepta \texttt{client\_request\_id} opcional; idempotente. \\
GET    & \http{/workouts/\{id\}}      & Detalle. \\
PUT    & \http{/workouts/\{id\}}      & \texttt{If-Match: <version>}. 409 si está stale. \\
DELETE & \http{/workouts/\{id\}}      & Hard delete con cascada a sets. \\
GET    & \http{/exercises}            & Catálogo, i18n por \texttt{Accept-Language}. \\
POST   & \http{/exercises}            & Ejercicio custom del usuario. \\
GET    & \http{/exercises/\{id\}/last-set} & Última serie registrada (autopopulado del cliente). \\
GET    & \http{/fatigue/summary}      & Llama al compute engine; cachea por usuario/día. \\
GET    & \http{/fatigue/weekly}       & Volumen semanal por grupo muscular (12 grupos). \\
GET    & \http{/routines}             & Catálogo, opcional \texttt{?level=}. \\
GET    & \http{/routines/\{id\}}      & Detalle con días/ejercicios. \\
\bottomrule
\caption{Endpoints REST.}\\
\end{longtable}

\subsubsection{Ejemplo: \texttt{POST /workouts} con idempotencia}

\begin{lstlisting}[language=bash,caption={Request idempotente con \texttt{client\_request\_id}.}]
POST /workouts
Authorization: Bearer <jwt>
Content-Type: application/json
X-Request-Id: 4e2b1d4f-...

{
  "name": "Push Day",
  "date": "2026-05-12T18:00:00Z",
  "duration_seconds": 3600,
  "client_request_id": "c5a30b7e-1d4f-4d7a-9b1f-2a2c3a8d9e10",
  "notes": null,
  "exercises": [
    {
      "exercise_id": "11111111-1111-1111-1111-111111111111",
      "order": 0,
      "sets": [
        {"weight": 60.0, "reps": 8, "rpe": 7.5, "order": 0},
        {"weight": 65.0, "reps": 5, "rpe": 8.5, "order": 1}
      ]
    }
  ]
}
\end{lstlisting}

\begin{lstlisting}[language=bash,caption={Respuesta 201 (primera vez) o 200 (retransmisión con la misma clave).}]
HTTP/1.1 201 Created
X-Request-Id: 4e2b1d4f-...
Content-Type: application/json

{
  "id": "f980dd08-...-b12d",
  "name": "Push Day",
  "date": "2026-05-12T18:00:00Z",
  "duration_seconds": 3600,
  "notes": null,
  "client_request_id": "c5a30b7e-...-9e10",
  "version": 0,
  "exercises": [ ... ]
}
\end{lstlisting}

\subsubsection{Ejemplo: \texttt{PUT /workouts/\{id\}} con If-Match}

\begin{lstlisting}[caption={Update condicional. El servidor compara \texttt{If-Match} con la \texttt{version} actual.}]
PUT /workouts/f980dd08-...-b12d
Authorization: Bearer <jwt>
If-Match: 1

{ "name": "Push Day (corregido)", "notes": "Mejor postura" }
\end{lstlisting}

\paragraph{Errores posibles.}
\texttt{200 OK} con \texttt{version} incrementada;
\texttt{404} si el workout no es del usuario;
\texttt{409 Conflict} con cuerpo \verb|{"error": "<localized>"}| si \texttt{If-Match} no coincide.

\subsubsection{Formato de errores}

Todos los errores siguen el mismo contrato:

\begin{lstlisting}
{ "error": "Workout has been modified by another session ... " }
\end{lstlisting}

Para validaciones de campos, el handler agrega \texttt{details}:

\begin{lstlisting}
{
  "error": "Validation failed",
  "details": ["email: must be a well-formed email address",
              "password: size must be between 8 and 255"]
}
\end{lstlisting}

\subsection{Seguridad}

\subsubsection{JWT RS256}

Los tokens se firman con RSA 2048 bits. El par de claves se carga desde variables de entorno (\texttt{JWT\_PRIVATE\_KEY}, \texttt{JWT\_PUBLIC\_KEY}) en formato PKCS8/X.509 base64. No existe ningún fallback en \texttt{application.yml}; si las variables faltan, el bean \cls{JwtTokenProvider} falla en el arranque.

\begin{table}[H]
\centering
\begin{tabular}{lp{0.5\textwidth}}
\toprule
\textbf{Claim} & \textbf{Significado} \\
\midrule
\texttt{sub}  & UUID del usuario. Lo escribe \cls{RbacFilter} en el \texttt{Authentication.name}. \\
\texttt{iss}  & \texttt{strengthlabs} \\
\texttt{iat}  & Fecha de emisión. \\
\texttt{exp}  & Expiración (24\,h por defecto). \\
\texttt{jti}  & UUID único usado como identificador para revocación. \\
\texttt{role} & \texttt{USER} o \texttt{ADMIN}. \\
\texttt{typ}  & \texttt{access} o \texttt{refresh}. \\
\bottomrule
\end{tabular}
\caption{Claims del JWT.}
\end{table}

\subsubsection{Refresh tokens y revocación}

El \texttt{POST /auth/refresh} consume el token actual y emite un par nuevo. El JTI anterior se añade a \cls{RevokedTokenStore} hasta su expiración natural. Este almacén es un \texttt{ConcurrentHashMap} en memoria — adecuado para el alcance académico, no para producción multi-instancia. La sección \textbf{Limitaciones} discute la migración a Redis.

\subsubsection{Rate limiting}

\cls{LoginRateLimitFilter} cuenta intentos por IP en una ventana deslizante (5 intentos / 15 min por defecto, configurable). Responde \texttt{429} con cabecera \texttt{Retry-After} cuando se excede. Aplica a \texttt{/auth/login}, \texttt{/auth/register}, \texttt{/auth/refresh}.

\subsubsection{Cabeceras y CORS}

\cls{SecurityConfig} configura:
\begin{itemize}[leftmargin=*]
  \item \texttt{Content-Type-Options: nosniff}
  \item \texttt{X-Frame-Options: deny}
  \item \texttt{Referrer-Policy: no-referrer}
  \item HSTS condicional (\texttt{security.hsts.enabled}) con \texttt{includeSubDomains} y 1 año.
  \item CORS configurable vía \texttt{security.cors.allowed-origins} (por defecto \texttt{localhost:*} y \texttt{10.0.2.2:*}).
\end{itemize}

\subsection{Idempotencia y concurrencia}

\subsubsection{Idempotencia de POST /workouts}

Cuando el cliente envía \texttt{client\_request\_id}, el controlador hace dos chequeos en orden:

\begin{enumerate}[leftmargin=*]
  \item Consulta \texttt{findByUserIdAndClientRequestId}. Si existe, devuelve esa entidad con \texttt{200 OK}.
  \item Si no, intenta insertar. El índice único parcial \texttt{(user\_id, client\_request\_id) WHERE client\_request\_id IS NOT NULL} es la red de seguridad: ante dos requests racing, sólo uno tendrá éxito y el otro recibe \texttt{DataIntegrityViolationException}, mapeada a \texttt{409}.
\end{enumerate}

\subsubsection{Optimistic locking de PUT /workouts/\{id\}}

\cls{WorkoutJpaEntity} lleva \texttt{@Version Long version}. Hibernate añade automáticamente la cláusula \texttt{WHERE version = ?} y compara filas afectadas en cada UPDATE. Además, el controlador inspecciona la cabecera \texttt{If-Match}; si el cliente envía una versión obsoleta, devuelve \texttt{409} de inmediato sin tocar la BD.

Importante: el método PUT usa \texttt{saveAndFlush} (no \texttt{save}) porque \texttt{@Version} se incrementa al \textit{flush}. Si la respuesta serializa antes del flush, el cliente recibe la versión vieja y queda atascado en \texttt{409} para siempre. Esto se detectó durante el primer despliegue en producción y está cubierto por una aserción explícita en \cls{WorkoutControllerIT.putRejectsStaleVersion}.

\subsubsection{Cache diario de fatiga}

\cls{GetFatigueSummaryUseCase.persistMetrics} verifica existencia y captura \texttt{DataIntegrityViolationException} para tolerar dos requests concurrentes del mismo usuario en el mismo día. Sin esto, el unique constraint \texttt{(user\_id, metric\_date)} levantaba excepción ruidosa en logs.

\subsection{Observabilidad}

\cls{CorrelationIdFilter} (\texttt{HIGHEST\_PRECEDENCE}) extrae o genera \texttt{X-Request-Id}, lo introduce en el MDC (\texttt{requestId}) y lo echoes en la respuesta. SLF4J/Logback están configurados para incluir \texttt{[requestId]} en el formato. Esto permite:

\begin{itemize}[leftmargin=*]
  \item Cruzar logs de servidor con el id que el cliente Dio inyecta en cada request.
  \item Trazar reintentos del SyncManager: como reintenta con el mismo \texttt{client\_request\_id}, basta filtrar por ese campo en la BD para ver toda la historia.
\end{itemize}

% =============================================================================
\section{Compute Engine: Python FastAPI}
% =============================================================================

\subsection{Endpoints}

\begin{table}[H]
\centering
\begin{tabularx}{\textwidth}{lXX}
\toprule
\textbf{Método} & \textbf{Ruta} & \textbf{Descripción} \\
\midrule
POST & \http{/compute/fatigue} & Recibe lista de sesiones \texttt{\{user\_id, date, duration\_minutes, rpe\}}; devuelve ATL/CTL/TSB, monotony, strain, ramp\_rate, readiness y serie ATL. \\
POST & \http{/compute/risk}    & Recibe ACWR, TSB, ramp\_rate, monotony; devuelve injury\_risk, overtraining\_risk, composite y \texttt{risk\_level} (low/moderate/high/critical), plus recomendaciones. \\
POST & \http{/plan}            & (Out of scope para el cliente actual.) \\
\bottomrule
\end{tabularx}
\caption{Superficie del microservicio Python.}
\end{table}

\subsection{Modelo de fatiga}

ATL y CTL se calculan con un promedio exponencial:

\[
ATL_t = \alpha_a \cdot load_t + (1 - \alpha_a) \cdot ATL_{t-1}, \quad \alpha_a = \frac{2}{7+1}
\]
\[
CTL_t = \alpha_c \cdot load_t + (1 - \alpha_c) \cdot CTL_{t-1}, \quad \alpha_c = \frac{2}{42+1}
\]
\[
TSB_t = CTL_t - ATL_t, \quad ACWR_t = ATL_t / \max(CTL_t, \epsilon)
\]

El \texttt{load} de cada sesión se aproxima como \texttt{duration\_minutes} $\times$ RPE (TRIMP simplificado). La serie ATL se devuelve completa para que el cliente pueda graficar tendencia sin volver a calcularla.

\subsection{Modelo de riesgo}

El servicio combina ACWR, TSB y ramp\_rate con pesos para producir un \texttt{composite\_risk\_score} y una clasificación discreta. Adicionalmente emite \texttt{risk\_flags} (códigos como \texttt{HIGH\_ACWR}, \texttt{DEEP\_FATIGUE}, \texttt{ELEVATED\_RAMP\_RATE}) que el backend traduce a strings localizados mediante \texttt{messages\_en.properties} / \texttt{messages\_es.properties}.

\subsection{Falla controlada}

\cls{PythonComputeAdapter} llama por HTTP. El \texttt{RestTemplate} está construido con \texttt{connectTimeout=2s} y \texttt{readTimeout=5s}. Cualquier excepción se atrapa y se devuelve \cls{FatigueComputeResult.unavailable()} con \texttt{computeAvailable=false}. El cliente Flutter conoce ese flag (\cls{FatigueSummary.hasComputeData}) y degrada la UI: muestra el volumen semanal calculado en backend pero oculta los gauges dependientes del compute.

% =============================================================================
\section{Frontend: Flutter}
% =============================================================================

\subsection{Estructura feature-first}

\begin{lstlisting}[caption={Estructura abreviada del módulo \texttt{lib/}.}]
lib/
|-- app.dart
|-- main.dart
|-- core/
|   |-- connectivity/        ConnectivityCubit
|   |-- constants/           api_constants.dart (baseUrl: api.blocsa.com)
|   |-- database/            Drift schema + app_database.g.dart
|   |-- demo/                DemoMode + mocks
|   |-- network/             DioClient + interceptores
|   |-- router/              GoRouter + guards
|   |-- storage/             TokenStorage (flutter_secure_storage)
|   '-- sync/                SyncManager
|-- features/
|   |-- auth/
|   |   |-- data/            AuthRepository, UserCache, SessionRestoreResult
|   |   '-- presentation/    AuthCubit + LoginPage/RegisterPage
|   |-- workouts/
|   |   |-- data/            WorkoutRepository, OfflineWorkoutService,
|   |   |                    ActiveWorkoutDraftStore
|   |   |-- domain/          Workout, Exercise, WorkoutSet, MuscleGroup
|   |   '-- presentation/    Cubits, pages
|   |-- fatigue/             repository + parser + dashboard
|   |-- routines/
|   |-- plan/
|   |-- export/              generacion local xlsx/csv
|   |-- settings/
|   '-- onboarding/
|-- shared/
|   |-- widgets/             OfflineBanner, etc.
|   '-- utils/
'-- l10n/                    ARB + generated
\end{lstlisting}

\subsection{Estado: Bloc / Cubit}

Cada feature relevante expone uno o más Cubits:
\cls{AuthCubit}, \cls{WorkoutsCubit}, \cls{ActiveWorkoutCubit}, \cls{FatigueCubit}, \cls{RoutinesCubit}, \cls{SettingsCubit}, \cls{ConnectivityCubit}. \texttt{setState} se permite únicamente para estado UI puramente local (visibilidad de un input, expansión de un menú).

\subsection{Capa de red}

\cls{DioClient} encadena tres interceptores en el orden:

\begin{enumerate}[leftmargin=*]
  \item \cls{CorrelationIdInterceptor}: agrega \texttt{X-Request-Id} (UUID v4) salvo que el llamante ya haya inyectado uno propio. El backend lo refleja en la respuesta y en cada línea de log.
  \item \cls{\_AuthInterceptor}:
    \begin{itemize}[leftmargin=*]
      \item Antes de cada request comprueba la expiración del access token (parseando el \texttt{exp} sin firmar otra vez) y refresca proactivamente.
      \item Si llega \texttt{401}, refresca y reintenta una sola vez.
      \item Refreshes concurrentes se colapsan mediante un \texttt{Completer<bool>} compartido para evitar tormenta de refreshes en paralelo.
    \end{itemize}
  \item \cls{RetryInterceptor}: reintenta sólo \texttt{GET} en errores transitorios (\texttt{connectionError}, \texttt{timeout}, 502, 503, 504) con backoff exponencial. POST/PUT/DELETE \textbf{nunca} se reintentan a este nivel — viajan por la cola de sync, que sí garantiza idempotencia.
\end{enumerate}

\subsection{Base de datos local: Drift}

\begin{lstlisting}[caption={Esquema Drift (Dart, fragmento).}]
@DataClassName('WorkoutRow')
class WorkoutsTable extends Table {
  TextColumn   get id              => text()();
  TextColumn   get serverId        => text().nullable()();
  TextColumn   get clientRequestId => text().unique()();
  TextColumn   get name            => text()();
  DateTimeColumn get date          => dateTime()();
  IntColumn    get durationSeconds => integer()();
  TextColumn   get notes           => text().nullable()();
  IntColumn    get version         => integer().withDefault(const Constant(0))();
  TextColumn   get exercisesJson   => text().withDefault(const Constant('[]'))();
  TextColumn   get syncState       =>
      text().withDefault(const Constant('synced'))();
  DateTimeColumn get localUpdatedAt  => dateTime()();
  DateTimeColumn get serverUpdatedAt => dateTime().nullable()();
  @override Set<Column> get primaryKey => {id};
}
\end{lstlisting}

El esquema incluye además \texttt{SyncQueueTable} (cola de operaciones con \texttt{attempts}, \texttt{lastError}, \texttt{nextRetryAt}) y \texttt{CachedUserTable} (un único registro con el usuario autenticado para el flujo offline de auth).

\subsection{Sesión activa persistente}

\cls{ActiveWorkoutCubit} serializa cada \texttt{emit} a \texttt{shared\_preferences} vía \cls{ActiveWorkoutDraftStore}. Si el SO mata la app a mitad de sesión, el siguiente arranque ofrece recuperar el draft. Para evitar que escrituras rápidas en serie produzcan que un save antiguo sobreescriba al más reciente, los saves se serializan a través de un único \texttt{Future} chain (\texttt{\_saveQueue}).

\subsection{Internacionalización}

ARB + \texttt{flutter\_localizations}. Dos locales: \texttt{en} y \texttt{es}. El cliente envía \texttt{Accept-Language} a través de \cls{DioClient.languageCode}, y el backend resuelve mensajes y nombres de ejercicios mediante \texttt{MessageSource} + columnas \texttt{name\_es} del catálogo.

% =============================================================================
\section{Estrategia offline-first}
% =============================================================================

\subsection{Principio rector}

\textit{El cliente nunca depende del servidor para confirmar una mutación.} La sesión es válida si hay token; los datos vienen de Drift; el servidor es un \textit{write-behind} alimentado por la cola.

\subsection{Máquina de estados del workout local}

\begin{figure}[H]
\centering
\begin{tikzpicture}[
  state/.style={rectangle,draw,rounded corners,minimum height=0.9cm,minimum width=2.3cm,align=center,font=\small},
  arrow/.style={-Latex,thick,font=\scriptsize},
]
  \node[state,fill=blue!10] (pc)   {pendingCreate};
  \node[state,fill=green!12,right=3.4cm of pc] (sync) {synced};
  \node[state,fill=yellow!15,below=1.4cm of sync] (pu) {pendingUpdate};
  \node[state,fill=red!10,below=1.4cm of pc]    (pd) {pendingDelete};
  \node[state,fill=gray!10,right=3.4cm of sync] (gone) {(eliminado)};

  \draw[arrow] (pc) -- node[above]{POST OK} (sync);
  \draw[arrow] (sync) -- node[right]{editar} (pu);
  \draw[arrow] (pu.west) to[bend left=25] node[left]{PUT OK} (sync.south west);
  \draw[arrow] (sync) -- node[below right,yshift=2pt]{borrar} (pd);
  \draw[arrow] (pd) -- node[below]{DELETE OK} (gone);
  \draw[arrow,dashed] (pc) to[bend right=20] node[left]{cancelar} (pd);
\end{tikzpicture}
\caption{Máquina de estados del campo \texttt{syncState} de una fila de \texttt{workouts}.}
\end{figure}

\subsection{Flujo end-to-end: crear workout sin red}

\begin{figure}[H]
\centering
\begin{tikzpicture}[
  actor/.style={font=\small\bfseries},
  msg/.style={-Latex,thick,font=\scriptsize},
  msgr/.style={Latex-,thick,font=\scriptsize},
  every text node part/.style={align=left},
]
  \node[actor] (u)   at (0,0)    {Usuario};
  \node[actor] (ui)  at (3,0)    {UI / Cubit};
  \node[actor] (svc) at (6.5,0)  {OfflineService};
  \node[actor] (db)  at (10,0)   {Drift};
  \node[actor] (sm)  at (13.5,0) {SyncManager};
  \node[actor] (api) at (17,0)   {Backend};

  \foreach \x in {0,3,6.5,10,13.5,17}
    \draw[dashed,gray] (\x,-0.2) -- (\x,-7.5);

  \draw[msg] (0,-0.7) -- node[above]{Tap Finish} (3,-0.7);
  \draw[msg] (3,-1.3) -- node[above]{createWorkout(draft)} (6.5,-1.3);
  \draw[msg] (6.5,-1.9) -- node[above]{INSERT pendingCreate} (10,-1.9);
  \draw[msgr] (6.5,-2.5) -- node[below]{ok} (10,-2.5);
  \draw[msg] (6.5,-3.1) -- node[above]{enqueue(createWorkout)} (13.5,-3.1);
  \draw[msgr] (3,-3.7) -- node[below]{LocalWorkoutHandle} (6.5,-3.7);
  \draw[msg]  (0,-4.3) -- node[above]{(UI muestra workout)} (3,-4.3);
  \draw[msg]  (13.5,-5.0) -- node[above]{POST /workouts \{client\_request\_id\}} (17,-5.0);
  \draw[msgr] (13.5,-5.6) -- node[below]{201 \{id, version\}} (17,-5.6);
  \draw[msg]  (13.5,-6.2) -- node[above]{UPDATE syncState=synced} (10,-6.2);
\end{tikzpicture}
\caption{Secuencia simplificada de \texttt{createWorkout} offline-first.}
\end{figure}

\subsection{Reintentos idempotentes}

Cuando el primer POST falla por red, \cls{SyncManager} marca \texttt{attempts++}, calcula \texttt{nextRetryAt = now + backoff(attempts)} y guarda el mismo payload. La siguiente vuelta del bucle reusa el mismo \texttt{client\_request\_id}; si la primera transacción del servidor sí había completado y sólo la respuesta se perdió, el backend devuelve \texttt{200} con la entidad existente y la cola colapsa la operación.

\subsection{Backoff}

\[
\text{delay}(n) = \min(\text{maxBackoff},\, \text{initialBackoff} \cdot 2^{\,n-1}) \cdot (1 + \mathcal{U}(-0.25,\,0.25))
\]

con \texttt{initialBackoff=1s} y \texttt{maxBackoff=1h} por defecto. El jitter del 25\% evita reintentos sincronizados tras una caída masiva.

\subsection{Resolución de conflictos}

El alcance académico cubre el caso típico (un usuario, un dispositivo). En multi-dispositivo se aplica \textit{last-write-wins} controlado por la \texttt{version} del servidor: si el PUT recibe \texttt{409}, el cliente refresca la fila desde el servidor y descarta la copia local con notificación al usuario.

% =============================================================================
\section{Flujos clave (secuencias)}
% =============================================================================

\subsection{Registro y obtención de tokens}

\begin{figure}[H]
\centering
\begin{tikzpicture}[
  actor/.style={font=\small\bfseries},
  msg/.style={-Latex,thick,font=\scriptsize},
  msgr/.style={Latex-,thick,font=\scriptsize},
]
  \node[actor] (ui)  at (0,0)    {Flutter};
  \node[actor] (api) at (5,0)    {Backend};
  \node[actor] (db)  at (10,0)   {Postgres};
  \node[actor] (cache) at (14.5,0) {UserCache};

  \foreach \x in {0,5,10,14.5} \draw[dashed,gray] (\x,-0.2) -- (\x,-6.0);

  \draw[msg]  (0,-0.7) -- node[above]{POST /auth/register} (5,-0.7);
  \draw[msg]  (5,-1.3) -- node[above]{INSERT user (bcrypt)} (10,-1.3);
  \draw[msgr] (5,-1.9) -- node[below]{ok} (10,-1.9);
  \draw[msgr] (0,-2.5) -- node[below]{200 \{access\_token, refresh\_token\}} (5,-2.5);
  \draw[msg]  (0,-3.1) -- node[above]{GET /auth/me} (5,-3.1);
  \draw[msgr] (0,-3.7) -- node[below]{200 \{id, name, email\}} (5,-3.7);
  \draw[msg]  (0,-4.3) -- node[above]{save user} (14.5,-4.3);
\end{tikzpicture}
\caption{Registro y persistencia del usuario en cache local.}
\end{figure}

\subsection{Restauración de sesión}

\cls{AuthRepository.restoreSession} devuelve un valor de tipo \cls{SessionRestoreResult}:

\begin{table}[H]
\centering
\begin{tabularx}{\textwidth}{lX}
\toprule
\textbf{Resultado} & \textbf{Cuándo} \\
\midrule
\cls{SessionMissing}        & No hay tokens almacenados. \\
\cls{SessionRestored}       & \texttt{/auth/me} confirmó al usuario. Cache actualizado. \\
\cls{SessionRestored(fromCache: true)} & La red falló pero había un \cls{UserCache} válido. \\
\cls{SessionExpired}        & El servidor devolvió 401. Tokens y cache borrados. \\
\cls{SessionOfflineNoCache} & Sin red y sin cache: arranque inicial sin internet. \\
\bottomrule
\end{tabularx}
\caption{Variantes de \texttt{SessionRestoreResult}.}
\end{table}

\subsection{Cálculo de fatiga con compute engine caído}

\begin{enumerate}[leftmargin=*]
  \item Cliente: \texttt{GET /fatigue/summary}.
  \item Backend: comprueba cache diaria; si existe, la devuelve.
  \item Cache miss: llama al microservicio Python con timeout 5\,s.
  \item Si el servicio responde, computa, persiste cache (con doble-check + catch de \texttt{DataIntegrityViolationException}) y devuelve.
  \item Si falla, devuelve \texttt{FatigueComputeResult.unavailable()} con valores neutros y \texttt{compute\_available=false}; \textbf{no cachea} para permitir reintento más tarde.
  \item Cliente: detecta \texttt{compute\_available=false} y oculta los gauges dependientes.
\end{enumerate}

% =============================================================================
\section{Despliegue}
% =============================================================================

\subsection{Topología actual}

\begin{itemize}[leftmargin=*]
  \item Host único Debian/Ubuntu.
  \item Postgres y Redis en systemd (\textit{native}, no contenedor).
  \item \texttt{strengthlabs.service}: ejecuta el JAR Spring Boot en \texttt{:8080}.
  \item \texttt{strengthlabs-compute.service}: uvicorn FastAPI en \texttt{:8001}.
  \item \texttt{cloudflared.service}: \textit{named tunnel} que mapea \texttt{api.blocsa.com} a \texttt{localhost:8080}. Sobrevive reboots, no rota credenciales, ligado a la cuenta.
\end{itemize}

\subsection{Variables de entorno requeridas}

\begin{table}[H]
\centering
\begin{tabularx}{\textwidth}{lX}
\toprule
\textbf{Variable} & \textbf{Uso} \\
\midrule
\texttt{DB\_PASSWORD}         & Password del rol \texttt{stlabs}. \\
\texttt{JWT\_PRIVATE\_KEY}    & PKCS8 base64 (sin headers). \\
\texttt{JWT\_PUBLIC\_KEY}     & X.509 base64. \\
\texttt{GOOGLE\_CLIENT\_ID}   & Para verificar ID tokens de Google. \\
\texttt{COMPUTE\_ENGINE\_URL} & Típicamente \texttt{http://localhost:8001}. \\
\texttt{CORS\_ALLOWED\_ORIGINS} & CSV de orígenes permitidos. \\
\texttt{PORT}                 & Puerto de escucha del backend. En el despliegue actual es \texttt{8080}. \\
\bottomrule
\end{tabularx}
\caption{Variables de entorno para producción.}
\end{table}

\subsection{Despliegue inicial sobre esquema previo}

El procedimiento documentado en \texttt{SERVER\_DEPLOY.md} \S6 cubre el caso del primer despliegue cuando ya existía un esquema construido por Hibernate sin Flyway: se baseliza V4 y se dejan correr V5 y V6 encima. El runbook contiene además un script de \textit{smoke tests} (\S5) que verifica registro, idempotencia, paginación, optimistic locking y echo de \texttt{X-Request-Id}.

\subsection{Configuración del cliente para producción}

El default de \cls{ApiConstants.baseUrl} es \texttt{https://api.blocsa.com}, de modo que \texttt{flutter run} en un dispositivo físico funciona sin sobreescrituras. Para desarrollo local se pasa \texttt{--dart-define=API\_BASE\_URL=...}.

% =============================================================================
\section{Estrategia de pruebas}
% =============================================================================

\subsection{Frontend (Flutter)}

\begin{itemize}[leftmargin=*]
  \item Parsing snake\_case ($\rightarrow$ Workout/Exercise/WorkoutSet), forzando UTC para fechas.
  \item \cls{FatigueResponseParser} defensivo: tolera \texttt{risk\_flags} como lista de strings o lista de mapas.
  \item \cls{AuthRepository.restoreSession} para los cinco resultados (online/offline/expirado/etc.).
  \item \cls{ActiveWorkoutCubit} con persistencia y recuperación, incluyendo emits en ráfaga.
  \item \cls{CorrelationIdInterceptor}: no pisa ids ya provistos, genera únicos.
  \item \cls{RetryInterceptor}: reintenta sólo GET en transitorios, abandona tras maxRetries.
  \item Drift schema: inserciones, constraints únicas, watch streams.
  \item \cls{SyncManager}: FIFO, backoff, abandono ante 4xx permanente.
  \item \cls{OfflineWorkoutService} end-to-end: crear offline $\rightarrow$ drain $\rightarrow$ \texttt{synced}, incluyendo reintento con mismo \texttt{client\_request\_id}.
\end{itemize}

\textbf{Ejecución:} \texttt{flutter test --concurrency=1} (las isolates compartiendo SQLite en memoria producen flakiness).

\subsection{Backend (JUnit 5 + Testcontainers)}

\begin{itemize}[leftmargin=*]
  \item \cls{AuthControllerIT}: registro, login, refresh, Google, logout.
  \item \cls{WorkoutControllerIT}: aislamiento por usuario, idempotencia, paginación, optimistic locking, version echo en el body de PUT.
  \item \cls{FatigueControllerIT}: caching y fallback ante compute caído.
  \item \cls{RateLimitIT}: 429 + \texttt{Retry-After}.
  \item \cls{JwtTokenProviderTest}, \cls{RevokedTokenStoreTest}, \cls{GetFatigueSummaryUseCaseTest}.
\end{itemize}

\textbf{Ejecución:} \texttt{JAVA\_HOME=/usr/lib/jvm/java-21-openjdk ./mvnw test}.

% =============================================================================
\section{Decisiones arquitectónicas (ADR)}
% =============================================================================

\subsection{ADR-001 — JWT RS256 en lugar de HS256}

\textbf{Contexto.} Se necesita firmar tokens emitidos por el backend y validados por el mismo servicio (y, ocasionalmente, por terceros vía JWKS).

\textbf{Decisión.} Usar RSA-256 (asimétrico). El backend firma con la clave privada; cualquier verificador externo descarga la pública en \texttt{/.well-known/jwks.json}.

\textbf{Consecuencias.} Operacionalmente más caro que HMAC (claves más grandes, generación de pares). Beneficio: los verificadores externos no necesitan compartir secretos, y rotar claves no obliga a revocar todos los tokens.

\subsection{ADR-002 — Offline-first con Drift en lugar de Hive}

\textbf{Contexto.} El cliente debe poder crear y leer entrenamientos sin red, y reconciliar al volver online.

\textbf{Decisión.} Drift (SQLite con codegen). Alternativas evaluadas: Hive, Isar.

\textbf{Razones.}
\begin{itemize}[leftmargin=*]
  \item Hive está estancado desde 2023; v2 sin mantenimiento activo.
  \item Isar v3 también detenido, v4 fue reescritura sin garantías.
  \item Drift es el ORM dart/Flutter más activo, con migraciones, streams reactivos y SQL real.
\end{itemize}

\textbf{Consecuencias.} Codegen requerido (\texttt{build\_runner build}). El .g.dart se commitea para que CI no necesite ejecutar build\_runner. SQLite real implica disponer de \texttt{libsqlite3} en el dispositivo (provisto por \texttt{sqlite3\_flutter\_libs}).

\subsection{ADR-003 — \texttt{client\_request\_id} en el body en lugar de cabecera \texttt{Idempotency-Key}}

\textbf{Contexto.} Necesitamos idempotencia en \texttt{POST /workouts}.

\textbf{Decisión.} Aceptar \texttt{client\_request\_id} como campo del cuerpo, y persistirlo como columna del workout (UUID nullable) con índice único parcial sobre \texttt{(user\_id, client\_request\_id)}.

\textbf{Alternativa.} Cabecera estándar \texttt{Idempotency-Key} + tabla auxiliar de keys.

\textbf{Razones.}
\begin{itemize}[leftmargin=*]
  \item La idempotencia es propia del recurso \texttt{workout}; modelarla como columna lo deja explícito en el esquema, sin tabla aparte.
  \item La columna permite consultas (``muéstrame todos los workouts que sincronizaron con la clave X'').
  \item La cabecera estándar habría requerido una tabla de ledger de keys con TTL, más complejidad operativa.
\end{itemize}

\textbf{Consecuencias.} El campo viaja en el JSON, no en headers; herramientas que esperan el estándar HTTP no lo verán automáticamente.

\subsection{ADR-004 — Spring Boot 4 en lugar de SB 3.x}

\textbf{Contexto.} El proyecto comenzó con SB 4.0.3.

\textbf{Decisión.} Mantenerse en SB4.

\textbf{Trampa.} Spring Boot 4 separó la autoconfiguración de Flyway en un módulo aparte (\texttt{spring-boot-flyway}). Sin esa dependencia, las migraciones se saltan silenciosamente. Esto causó el primer crash de despliegue (Hibernate corrió antes que Flyway). Documentado en \texttt{SERVER\_DEPLOY.md}.

\subsection{ADR-005 — Compute engine separado en Python}

\textbf{Contexto.} Los modelos de fatiga y riesgo son cálculo numérico iterativo. Implementarlos en Java requería traducir lógica que el dominio del fitness suele tener en Python.

\textbf{Decisión.} Servicio aparte en FastAPI, llamado por HTTP.

\textbf{Consecuencias.} Despliegue tiene dos procesos. El backend debe tolerar caídas del compute (timeout + fallback). El cliente debe distinguir ``valores cero porque no entrené'' de ``valores cero porque el compute está caído'': se usa el flag \texttt{compute\_available}.

% =============================================================================
\section{Limitaciones conocidas y trabajo futuro}
% =============================================================================

\begin{itemize}[leftmargin=*]
  \item \textbf{Revocación de tokens en memoria.} \cls{RevokedTokenStore} es un mapa local; un reinicio del proceso pierde la blacklist y, en una hipotética instalación multi-instancia, no se sincroniza. Migración futura a Redis.
  \item \textbf{Sin recuperación de contraseña.} El flujo de reset por email no está implementado.
  \item \textbf{Pull bidireccional manual.} El cliente no escucha cambios proactivos del servidor; sincroniza on-demand al abrir pantallas relevantes. WebSocket o long-poll quedan fuera de alcance.
  \item \textbf{Sin telemetría remota.} No hay Sentry ni Crashlytics. Los logs viven en el host.
  \item \textbf{Sin rate limiting distribuido.} \cls{LoginRateLimitFilter} cuenta intentos en memoria, por proceso.
  \item \textbf{Despliegue mono-host.} No hay alta disponibilidad. Postgres es local; copias de seguridad son responsabilidad del operador.
  \item \textbf{Exportes locales únicamente.} \texttt{.xlsx} y \texttt{.csv} se generan en el dispositivo con \texttt{excel} y \texttt{share\_plus}. No hay endpoint de descarga en el servidor (el antiguo \cls{ExportController} fue retirado).
\end{itemize}

% =============================================================================
\section{Referencias}
% =============================================================================

\begin{itemize}[leftmargin=*]
  \item Banister, E. (1991). \textit{Modeling elite athletic performance}. Bases del modelo ATL/CTL/TSB.
  \item Gabbett, T. (2016). \textit{The training–injury prevention paradox}. ACWR como predictor de lesión.
  \item Spring Boot 4 release notes — \url{https://github.com/spring-projects/spring-boot/wiki/Spring-Boot-4.0-Release-Notes}.
  \item Drift documentation — \url{https://drift.simonbinder.eu/}.
  \item Cloudflare Tunnel docs — \url{https://developers.cloudflare.com/cloudflare-one/connections/connect-networks/}.
  \item Repositorios del proyecto — \url{https://github.com/YeisenK/StrengthLabs} y \url{https://github.com/YeisenK/StrengthLabsBackend}.
\end{itemize}

\end{document}
